\documentclass[notes,10pt, aspectratio=169]{beamer}

\usepackage{pgfpages}
% These slides also contain speaker notes. You can print just the slides,
% just the notes, or both, depending on the setting below. Comment out the want
% you want.
% Note visibility is selected after package loading at the end of this preamble.
%\setbeameroption{show only notes} % Only notes
%\setbeameroption{show notes on second screen=right} % Both

\usepackage{helvet}
\usepackage[default]{lato}
\usepackage{array}
\usepackage{eurosym}
\usepackage{setspace}
\usepackage{animate}
\usepackage{multicol}
\usepackage{breqn}
\usepackage{subcaption}
\usepackage{caption}
\captionsetup{font=scriptsize, justification=raggedright,singlelinecheck=false}
\usepackage[export]{adjustbox}
\usepackage[scale=2]{ccicons}
\usepackage{tikz}
\newenvironment{wideitemize}{\itemize\addtolength{\itemsep}{10pt}}{\enditemize}
\usepackage{natbib}

\usepackage{tikz}
\usepackage{verbatim}
\setbeamertemplate{note page}{\pagecolor{yellow!5}\insertnote}
\usepackage{amsmath}
%\usepackage{mathpazo}
\usepackage{hyperref}
\usepackage{lipsum}
\usepackage{multimedia}
\usepackage{multirow}
\usepackage{graphicx}
\usepackage{dcolumn}
\usepackage{bbm}
\newcolumntype{d}[0]{D{.}{.}{5}}

\usepackage{changepage}
\usepackage{appendixnumberbeamer}
% \newcommand{\beginbackup}{
%   \newcounter{framenumbervorappendix}
%   \setcounter{framenumbervorappendix}{\value{framenumber}}
%   \setbeamertemplate{footline}
%   {
%      \leavevmode%
%      \hline
%      box{%
%       \begin{beamercolorbox}[wd=\paperwidth,ht=2.25ex,dp=1ex,right]{footlinecolor}%
% %         \insertframenumber  \hspace*{2ex} 
%       \end{beamercolorbox}}%
%      \vskip0pt%
%   }
%  }

 \setbeamertemplate{footline}{
    \hbox{%
    \begin{beamercolorbox}[wd=\paperwidth,ht=3ex,dp=1.5ex,leftskip=2ex,rightskip=2ex]{page footer}%
        %\usebeamerfont{title in head/foot}%
        %\insertshorttitle \hfill
        %\insertsection \hfill
        %\tableofcontents[currentsection] \hfill
        \insertsectionnavigationhorizontal{0.7\paperwidth}{}{} \hfill 
        \insertframenumber{} / \inserttotalframenumber
    \end{beamercolorbox}}%
}
\newcommand{\backupend}{
   \addtocounter{framenumbervorappendix}{-\value{framenumber}}
   \addtocounter{framenumber}{\value{framenumbervorappendix}} 
}


\usepackage{graphicx}
\usepackage[space]{grffile}
\usepackage{booktabs}

% These are my colors -- there are many like them, but these ones are mine.
\definecolor{blue}{RGB}{0,114,178}
\definecolor{red}{RGB}{213,94,0}
\definecolor{yellow}{RGB}{240,228,66}
\definecolor{green}{RGB}{0,158,115}

\setbeamercolor{section in head/foot}{fg=blue}

\useinnertheme[shadow]{rounded}
\setbeamercolor{block title example}{fg=black,bg=orange!50!white}
\setbeamercolor{block body example}{fg=black, bg=orange!30!white}


\hypersetup{
  colorlinks=false,
  linkbordercolor = {white},
  linkcolor = {blue}
}


%% I use a beige off white for my background
\definecolor{MyBackground}{RGB}{255,253,218}

%% Uncomment this if you want to change the background color to something else
%\setbeamercolor{background canvas}{bg=MyBackground}

%% Change the bg color to adjust your transition slide background color!
\newenvironment{transitionframe}{
  \setbeamercolor{background canvas}{bg=yellow}
  \begin{frame}}{
    \end{frame}
}

\setbeamercolor{frametitle}{fg=blue}
\setbeamercolor{title}{fg=black}
%\setbeamertemplate{footline}[frame number]
\setbeamertemplate{navigation symbols}{} 
\setbeamertemplate{itemize items}{-}
\setbeamercolor{itemize item}{fg=blue}
\setbeamercolor{itemize subitem}{fg=blue}
\setbeamercolor{enumerate item}{fg=blue}
\setbeamercolor{enumerate subitem}{fg=blue}
\setbeamercolor{button}{bg=white,fg=blue,}



% If you like road maps, rather than having clutter at the top, have a roadmap show up at the end of each section 
% (and after your introduction)
% Uncomment this is if you want the roadmap!
% \AtBeginSection[]
% {
%    \begin{frame}
%        \frametitle{Roadmap of Talk}
%        \tableofcontents[currentsection]
%    \end{frame}
% }
\setbeamercolor{section in toc}{fg=blue}
\setbeamercolor{subsection in toc}{fg=red}
%\setbeamersize{text margin left=1em,text margin right=1em} 
\setbeamersize{text margin left=0.7cm,text margin right=0.7cm} 


\title[]{\textcolor{blue}{Bank of Questions for Exams}}
\subtitle{Questions and worked solutions}
\author{25117 - Econometrics}
\institute[shortinst]{Universitat Pompeu Fabra}

%\date{September 25th, 2024}


% Each problem is followed by worked solution slides. Set \showsolutionsfalse
% below to compile a problem-only deck. Hidden \note entries retain instructor
% guidance. All data are hypothetical; no external assets are needed.
\newif\ifshowsolutions
\showsolutionstrue
\newcommand{\E}{\mathbb{E}}
\newcommand{\Var}{\operatorname{Var}}
\newcommand{\Cov}{\operatorname{Cov}}
\newcommand{\Corr}{\operatorname{Corr}}
\newcommand{\answerpart}[1]{\textcolor{blue}{\textbf{#1.}}\,}

\setbeameroption{hide notes} % Speaker notes hidden; worked solutions are visible.

\begin{document}

\begin{frame}[plain]
\maketitle
\end{frame}

% Source: Seminar 2, original Problem 1. Moved here for future exams.
\section{Sample dispersion}

\begin{frame}{Problem 1: How dispersed are five test scores?}
Five students obtain the following scores, measured in points:
\[
60, 60, 70, 80, 80.
\]
Treat the scores as i.i.d. observations from a population with finite variance.
\begin{enumerate}
\setlength{\itemsep}{9pt}
\item Find the sample mean and each student's deviation from it. Check the sum of the deviations.
\item Calculate the sample variance using the divisor $n-1$, and then the sample standard deviation. State their units.
\item What variance would the divisor $n$ give? Explain why $n-1$ is used when estimating the population variance, and identify the remaining degrees of freedom.
\end{enumerate}
\end{frame}
\note{All scores are hypothetical teaching data. Solution to Problem 1. The sample mean is 70; the deviations are -10,-10,0,10,10 and sum to zero. The sum of squared deviations is 400. The sample variance is 400/(5-1)=100 points squared and the sample standard deviation is sqrt(100)=10 points. Dividing by n instead gives 80 points squared. Estimating the mean imposes one linear constraint on the deviations, leaving n-1=4 degrees of freedom. Under the stated i.i.d. finite-variance assumptions, the expected sum of squared deviations is (n-1) times the population variance. Bessel's correction makes the sample variance unbiased; it does not generally make the sample standard deviation unbiased.}

\ifshowsolutions
% SOLUTION 1 START
\begin{frame}{Solution 1: Mean, deviations, and sample dispersion}
\small
\answerpart{1}The sample mean is
\[
\overline y=\frac1n\sum_{i=1}^n y_i
=\frac{60+60+70+80+80}{5}=\boxed{70\text{ points}}.
\]
\begin{center}
\begin{tabular}{c|rrrrr|r}
Student $i$ & 1 & 2 & 3 & 4 & 5 & Sum \\ \hline
$y_i-\overline y$ & $-10$ & $-10$ & 0 & 10 & 10 & 0 \\
$(y_i-\overline y)^2$ & 100 & 100 & 0 & 100 & 100 & 400
\end{tabular}
\end{center}
\answerpart{2}Using the sum of squared deviations,
\[
s_Y^2=\frac{\sum_{i=1}^n(y_i-\overline y)^2}{n-1}
=\frac{400}{5-1}=\boxed{100\text{ points}^2}.
\]
The standard deviation is the square root of the variance:
\[
s_Y=\sqrt{s_Y^2}=\sqrt{100}=\boxed{10\text{ points}}.
\]
\end{frame}

\begin{frame}{Solution 1: Why the divisor is $n-1$ (continued)}
\small
\answerpart{3}Dividing by $n$ would give
\[
\widetilde s_Y^2=\frac1n\sum_{i=1}^n(y_i-\overline y)^2
=\frac{400}{5}=\boxed{80\text{ points}^2},
\]
compared with $s_Y^2=100$ points squared using $n-1$.
\begin{itemize}
\setlength{\itemsep}{7pt}
\item Once the mean is fitted, $\sum_i(y_i-\overline y)=0$. Knowing four of the five deviations determines the fifth: there are $n-1=4$ degrees of freedom.
\item For i.i.d. observations with finite population variance $\sigma_Y^2$,
\[
\E\!\left[\sum_{i=1}^n(Y_i-\overline Y)^2\right]=(n-1)\sigma_Y^2.
\]
Consequently, dividing by $n$ has expectation $\frac{n-1}{n}\sigma_Y^2$, whereas dividing by $n-1$ has expectation $\sigma_Y^2$.
\end{itemize}
\begin{block}{Bessel's correction}
The divisor $n-1$ makes the \emph{variance estimator} unbiased. Its square root $s_Y$ is not generally an unbiased estimator of $\sigma_Y$.
\end{block}
\end{frame}
% SOLUTION 1 END
\fi

\end{document}
